
در نمودار، دو ناحیه‌ی جداگانه سایه‌خورده است.

ناحیه‌ی بالا-راست داخلِ $A$ و $C$ است ولی داخلِ $B$ نیست؛ پس:
\[
(A\cap C)\setminus B \,=\, A\cap C\cap B^{c}.
\]

ناحیه‌ی پایین-چپ داخلِ $A$ و $B$ است ولی داخلِ $C$ نیست؛ پس:
\[
(A\cap B)\setminus C \,=\, A\cap B\cap C^{c}.
\]

بنابراین کل ناحیه‌ی سایه‌خورده برابر اجتماعِ این دو قسمت است:
\[
\big((A\cap C)\setminus B\big)\;\cup\;\big((A\cap B)\setminus C\big).
\]

اگر از تعریف تفاضلِ متقارن یا همان \lr{symmetric difference} استفاده کنیم، داریم:
\[
\big((A\cap C)\setminus B\big)\cup\big((A\cap B)\setminus C\big)
\;=\;
A\cap\big((C\setminus B)\cup(B\setminus C)\big)
\;=\;
A\cap (B\,\triangle\,C).
\]
به بیان دیگر: نقاطی از $A$ که دقیقاً در یکی از $B$ یا $C$ قرار دارند.
